%%%%%%%%%%%%%%%%%%%%%%%%%%%%%%%%%%%%%%%%%%%%%%%%%%%%%%%%%%%%%%%%%%%%%%%%
%% Chang-Han Chen -- Curriculum Vitae
%% Clean single-column layout. Compile with pdflatex (run twice).
%%%%%%%%%%%%%%%%%%%%%%%%%%%%%%%%%%%%%%%%%%%%%%%%%%%%%%%%%%%%%%%%%%%%%%%%
\documentclass[11pt]{article}

\usepackage[a4paper,margin=18mm]{geometry}
\usepackage[T1]{fontenc}
\usepackage{lmodern}
\usepackage{enumitem}
\usepackage{titlesec}
\usepackage{xcolor}
\usepackage[hidelinks]{hyperref}

% ---- Colors -----------------------------------------------------------
\definecolor{accent}{HTML}{1A4E8A}   % section headings / links
\definecolor{rule}{HTML}{1A4E8A}
\definecolor{hl}{HTML}{B22222}        % subtle highlight (dates, key facts)

\hypersetup{colorlinks=true, urlcolor=accent, linkcolor=accent}

% ---- Layout -----------------------------------------------------------
\setlength{\parindent}{0pt}
\pagestyle{empty}

% Section headings: full-width, colored, bold, with a rule underneath.
\titleformat{\section}
  {\large\bfseries\color{accent}}{}{0pt}{}[\vspace{1pt}{\color{rule}\titlerule[1pt]}]
\titlespacing*{\section}{0pt}{12pt}{6pt}

% List styling: tight, no extra left indent.
\setlist[itemize]{leftmargin=1.4em, itemsep=2pt, topsep=2pt, parsep=0pt}

% ---- Helper macros ----------------------------------------------------
% \entry{title}{right-aligned date}
\newcommand{\entry}[2]{\textbf{#1}\hfill{\small #2}}
% \hl{text} -- subtle highlight
\newcommand{\hl}[1]{\textcolor{hl}{#1}}

\begin{document}

% ======================= Header =======================================
{\huge\bfseries Chang-Han Chen}\hfill{\small\color{gray}Last update: June 2026}

\vspace{4pt}
{\color{rule}\hrule height 1pt}
\vspace{6pt}

\begin{tabular}{@{}l@{}}
Tel: 608-886-6081 \quad\textbullet\quad
E-mail: \href{mailto:changhanc@berkeley.edu}{changhanc@berkeley.edu} \quad\textbullet\quad
Web: \href{https://chang-han-chen.github.io}{chang-han-chen.github.io}
\end{tabular}

% ======================= Education ====================================
\section{Education}
\begin{itemize}
  \item \entry{Ph.D. in Theoretical Physics, University of California, Berkeley}{expected \hl{end of 2026}}
  \item \entry{B.Sc. in Physics and Mathematics, Massachusetts Institute of Technology}{2022}\\
        GPA \hl{4.9/5}; Morse/Orloff award for best senior thesis.
\end{itemize}

% ======================= Skills =======================================
\section{Skills (besides physics)}
\begin{itemize}
  \item \textbf{Research:} von Neumann algebras, random matrix and free probability, language model training.
  \item \textbf{Coursework:} \href{https://ocw.mit.edu/courses/6-046j-design-and-analysis-of-algorithms-spring-2015/}{algorithm design},
        \href{https://www.argmin.net/p/convex-optimization-in-the-age-of}{convex optimization},
        probability theory, statistical learning theory.
\end{itemize}

% ======================= Research Interest ============================
\section{Research Interest}
\textit{What is the most puzzling, mysterious aspect of our universe?}

\vspace{4pt}
\textbf{Universality of holography: GR + QM = hologram?}
Instead of betting on one particular theory such as string theory or loop quantum gravity that are
``invisible'' to current technology, I believe in the universality of holography: \emph{any}
non-perturbative quantum theory consistent with general relativity is holographic, i.e.\ the
information in quantum gravity lives on the boundary. The goal of my Ph.D.\ is to explore as many
aspects of this universality as possible, e.g.\ two-dimensional toy models and von Neumann algebras.

\vspace{4pt}
For my recent interest in AI, see \href{https://chang-han-chen.github.io/blog/a-new-chapter/}{this blog}.

% ======================= Selected Projects ============================
\section{Selected Projects}
\begin{itemize}
  \item \entry{\href{https://github.com/Chang-Han-Chen/autoresearch-moe}{\textcolor{hl}{MoE autoresearch}}}{2026}\\
        Explored LLM-agent-driven ``autoresearch'' on autoregressive mixture-of-experts transformers,
        stacking interventions to a 1.75\% bits-per-byte gain that persists under compute-optimal scaling.
  \item \entry{\href{https://github.com/Chang-Han-Chen/sample-efficient-dlm}{\textcolor{black}{Optimal block-size curriculum for diffusion language models}}}{2026 Spring}\\
        Designed a block-size curriculum for pretraining block-diffusion language models, sweeping the
        autoregressive-vs-diffusion training budget and finding a consistent sample-efficiency optimum near 30\%.
  \item \entry{\href{https://arxiv.org/abs/2506.04311}{\textcolor{black}{An apologia for islands}}\ \normalfont(with S.\ Antonini, H.\ Maxfield, G.\ Penington)}{2025 Spring}\\
        Proved that entanglement islands and the Page curve emerge even with \emph{massless} gravitons and
        no external reservoir, using gauge-invariant gravitational operators to show black-hole interior
        operators are non-perturbatively encoded in Hawking radiation.
  \item \entry{\href{https://arxiv.org/abs/2406.02116}{\textcolor{hl}{A clock is just a way to tell the time}}\ \normalfont(with G.\ Penington)}{2024 Spring}\\
        Resolved prevalent misconceptions and pioneered practical applications of abstract mathematical
        tools to cosmological modeling, with work that became fairly influential in the field.
  \item \entry{\href{https://inspirehep.net/files/c1b52950194cd41895084217803489d7}{\textcolor{black}{Geometry and entanglement in AdS/CFT and beyond}}}{2022 Spring}\\
        Independent senior thesis, awarded MIT's Morse/Orloff prize for best senior thesis in theoretical physics.
\end{itemize}

% ======================= Honors and Awards ============================
\section{Honors and Awards}
\begin{itemize}
  \item \entry{\href{https://physics.mit.edu/academic-programs/student-awards/}{Morse/Orloff Research Award} -- best senior thesis in theoretical physics, MIT}{2022}
  \item \entry{Gold medal, \href{https://ipho2018.pt/content/exams}{International Physics Olympiad} (IPhO)}{2018}
\end{itemize}

% ======================= References ===================================
\section{References}
\begin{tabular}{@{}lll@{}}
Prof.\ \textbf{Hong Liu} & Prof.\ \textbf{Geoff Penington} & Prof.\ \textbf{Aron Wall}\\
MIT & UC Berkeley & U.\ of Cambridge\\
\href{mailto:hong_liu@mit.edu}{hong\_liu@mit.edu} &
\href{mailto:geoffp@berkeley.edu}{geoffp@berkeley.edu} &
\href{mailto:aroncwall@gmail.com}{aroncwall@gmail.com}\\
\end{tabular}

\end{document}
